% Figure 2: The tangent space of a submanifold
\begin{tikzpicture}[line cap=round,line join=round,>=Stealth,scale=0.9]
  % Left: curved sheet N containing the curve M through x
  \draw[thick] (-4.6,0.1) .. controls (-4.2,1.0) and (-3.6,1.6) .. (-2.9,1.75)
    .. controls (-2.4,1.85) and (-1.9,1.7) .. (-1.6,1.45)
    .. controls (-1.95,1.2) and (-2.2,0.8) .. (-2.25,0.35)
    .. controls (-2.9,0.35) and (-3.9,0.25) .. (-4.6,0.1);
  \node at (-1.62,1.85) {$N$};
  % the 1-dimensional submanifold M on the sheet
  \draw[very thick] (-4.25,0.45) .. controls (-3.65,0.75) and (-3.15,1.05) .. (-2.85,1.58);
  \node at (-2.72,1.95) {$M$};
  \fill (-3.45,0.92) circle (1.6pt);
  \node[below right] at (-3.47,0.9) {\footnotesize $x$};
  % Right: tangent plane TN_x with the line TM_x through 0
  \draw[thin] (0.4,0.2) -- (3.4,0.55) -- (4.3,1.75) -- (1.3,1.4) -- cycle;
  \node at (4.1,0.35) {$TN_x$};
  \fill (2.35,0.98) circle (1.6pt);
  \node[below right] at (2.37,0.96) {\footnotesize $0$};
  \draw[thick] (1.0,0.28) -- (3.7,1.68);
  \node at (4.0,1.95) {$TM_x$};
\end{tikzpicture}
